\documentclass{article}
\usepackage{amsmath, amssymb, amsfonts}
\title{The Block Theorem}
\begin{document}
\author{Amey Joshi}
\maketitle
If $\vec{R}$ is a translation vector in the primary lattice then the potential
function is periodic such that
\begin{equation}\label{e1}
U(\vec{r} + \vec{R}) = U(\vec{r}).
\end{equation}
Let $T_{\vec{R}}$ be a translation operator defined as
\begin{equation}\label{e2}
T_{\vec{R}} f(\vec{r}) = f(\vec{r} + \vec{R}).
\end{equation}
Then,
\[
[U, T_{\vec{R}}]\psi(\vec{r}) = U(\vec{r})\psi(\vec{r} + \vec{R}) - U(\vec{r} + 
\vec{R})\psi(\vec{r} + \vec{R}) = U(\vec{r})\psi(\vec{r}) - U(\vec{r})
\psi(\vec{r}) = 0.
\]
Likewise,
\[
[\nabla^2, T_{\vec{R}}]\psi(\vec{r}) = \nabla^2\psi(\vec{r} + \vec{R}) - 
T_{\vec{R}}\left(\nabla^2 \psi(\vec{r})\right).
\]
Since
\[
\frac{\partial f}{\partial x} = \frac{\partial f}{\partial (x+a)}
\frac{\partial (x+a)}{\partial x} = \frac{\partial f}{\partial (x+a)},
\]
we have $T_{\vec{R}}\left(\nabla^2 \psi(\vec{r})\right) = \nabla^2 \psi(\vec{r}
+ \vec{R})$ so that $[\nabla^2, T_{\vec{R}}] = 0$.

From these two commutation relations it follows that $T_{\vec{R}}$ commutes with
the Hamiltonian $H$ and they have simultaneous eigenstates. If the eigenvalue
equation of $T_{\vec{R}}$ is
\begin{equation}\label{e3}
T_{\vec{R}}\psi(\vec{r}) = C_{\vec{R}}\psi(\vec{r})
\end{equation}
If $\vec{R}_1$ and $\vec{R}_2$ are translation vectors then so is their sum.
Therefore, we must have
\begin{equation}\label{e4}
C_{\vec{R}_1 + \vec{R}_2} = C_{\vec{R}_1}C_{\vec{R}_2}.
\end{equation}
If $\psi(\vec{r})$ is normalised then so is $\psi(\vec{r} + \vec{R})$. As a 
result, we also have
\begin{equation}\label{e5}
|C_{\vec{R}}|^2 = 1.
\end{equation}
Equations \eqref{e4} and \eqref{e5} can be satisfied if
\begin{equation}\label{e6}
C_{\vec{R}} = e^{i\vec{k}\cdot\vec{R}}
\end{equation}
so that
\[
\psi(\vec{r} + \vec{R}) = \psi(\vec{r}) e^{i\vec{k}\cdot\vec{R}} \Rightarrow
\psi(\vec{r} + \vec{R})e^{-i\vec{k}\cdot\vec{R}} = \psi(\vec{r}) \Rightarrow
\psi(\vec{r} + \vec{R})e^{-i\vec{k}\cdot(\vec{r} + \vec{R})} = 
\psi(\vec{r})e^{-\vec{k}\cdot\vec{r}}.
\]
Thus, the function $\psi(\vec{r})e^{\vec{k}\cdot\vec{r}}$ has a period equal to
any lattice translation vector. We write the energy eigenstate as
\begin{equation}\label{e7}
\psi_{n\vec{k}} = \frac{1}{\sqrt{V}}u_{n\vec{k}} e^{i\vec{k}\cdot\vec{r}},
\end{equation}
where $n$ is the band-index and $V$ is the volume of the primitive cell. The
function $u_{n\vec{k}}$ is called the cell function and $\psi_{n\vec{k}}$ is
called the Block function.
\end{document}
